\chapter{Нейронні мережі, які вміють навчатися: додаємо \nt{DOPA}}
\chaptermark{Мережі, які навчаються}
\label{sec:dofa}
\begin{chaptergoals}
\item чому синапсу доступні лише локальні сигнали і чому це виключає зворотне поширення помилки;
\item чого навчає правило Гебба: напрямку, в якому входи змінюються найсильніше;
\item як слід і широкомовний сигнал $\delta$ перетворюють шум на кроки вгору градієнтом винагороди;
\item як критик дає $\delta$ у неперервному часі і чому один спільний провід навчає повільно.
\end{chaptergoals}

\section{Правила гри}

В усіх схемах попередніх розділів ваги ставив інженер. В організмі інженера немає. Ваги мають установлюватися самі, під час роботи, і так, щоб мережа робила щось корисне. Це й називають навчанням.

До попередніх правил додається ще одне. Синапс змінює свою вагу сам, і знати він може лише те, що відбувається \emph{поруч із ним}: активність відправника $x$, активність отримувача $y$ і ті речовини, що до нього доходять. Ніхто не може надіслати синапсу персональну поправку.

Це виключає головний метод штучних нейронних мереж --- зворотне поширення помилки. Там помилка виходу проходить мережу у зворотному напрямку тими самими вагами, в окрему фазу після прямого проходу, і кожна вага отримує свою поправку. Для цього потрібні зворотні проводи, що точно повторюють прямі, і спільний такт. Ні того, ні іншого в мозку не знайдено, принаймні в тій формі, що в штучних мережах~\cite{Lillicrap2020,WhittingtonBogacz2019}.

Зміну ваги записуватимемо як повільний неперервний процес:
\begin{equation}
\frac{\dd w}{\dd t} \;=\; \eta\,\Phi(\text{те, що відомо синапсу}),
\label{eq:plasticity}
\end{equation}
де $\eta$ --- швидкість навчання, настільки мала, що за час одного імпульсу вага майже не змінюється. Увесь розділ --- про те, якою може бути функція $\Phi$.

\section{Де зберігається вага}
\label{sec:weightstore}

Що таке синапс, який «змінює свою вагу»? Зв'язок має два боки: закінчення проводу відправника (\emph{пресинаптичний} бік) і ділянку мембрани отримувача з рецепторами (\emph{постсинаптичний} бік), а між ними --- вузька щілина. У \nt{GLUT}-зв'язків мембрана отримувача зазвичай розташована на \emph{шипику} --- маленькому виступі вхідної гілки отримувача. Вагу зв'язку зручно розкласти на три множники~\cite{delCastillo1954}:
\begin{empheq}[box=\keybox]{equation}
w \;=\; N_s\,p\,A_1 .
\label{eq:quantal}
\end{empheq}
Тут $N_s$ --- кількість місць викиду між двома нейронами, $p$ --- імовірність, що імпульс викине порцію медіатора в одному місці (та сама $p$ з розділу~\ref{sec:code}), $A_1$ --- відповідь отримувача на одну порцію, тобто скільки рецепторів її вловлює. Множник $p$ живе у відправника, $A_1$ --- в отримувача, а $N_s$ потребує обох боків: нові місця викиду виростають, старі зникають, і це триває все життя.

\emph{Вирішує отримувач.} У типового \nt{GLUT}-синапсу один із рецепторів глутамату проводить струм, лише якщо збіглися дві умови: глутамат прийшов від відправника, і мембрана отримувача вже збуджена~\cite{Nowak1984}. Це правило Гебба, вбудоване в одну молекулу, --- детектор збігів із розділу~\ref{sec:glutcomputer}, поставлений на вхід самого синапсу. Спрацювавши, рецептор впускає в отримувача кальцій, і кальцій запускає зміну ваги. Отримувач --- єдине місце, де видно водночас і вхід, і вихід, тому рішення ухвалюється в нього.

\emph{Виконати рішення може будь-який бік.} Найбільш вивчене довготривале посилення відбувається в отримувача: у мембрану вбудовуються нові рецептори~\cite{Malinow2002}, і шипик росте~\cite{Matsuzaki2004} --- росте $A_1$. Але отримувач може змінити і $p$: він викидає речовину, яка йде через щілину назад, до відправника, --- зворотний канал із додатка~\ref{sec:othermediators}, --- і той починає викидати менше медіатора~\cite{WilsonNicoll2001}. А короткочасні зміни ваги, виснаження і полегшення з розділу~\ref{sec:code}, --- це зміни $p$, які відправник веде сам за своєю недавньою активністю.

Для інженера регістр ваги розділено між двома мікросхемами. Правило навчання виконує отримувач, а записати результат він може і в себе, і у відправника, через зворотний канал. Тому слово «локально» в правилах цього розділу означає не один бік зв'язку, а весь зв'язок цілком.

\section{Навчатися без учителя}

Найпростіше, що може зробити синапс, --- посилитися, коли відправник і отримувач активні разом:
\begin{equation}
\frac{\dd w}{\dd t} \;=\; \eta\,x\,y .
\label{eq:hebb}
\end{equation}
Це правило Гебба~\cite{Hebb1949}. Воно локальне і не потребує годинника. Але в нього та сама біда, що в \nt{GLUT}-мережі в розділі~\ref{sec:glutcomputer}: додатний зворотний зв'язок. Сильна вага робить $y$ більшим, великий $y$ ще посилює вагу, і ваги ростуть без меж. Ліки --- від'ємний зворотний зв'язок за вагою: нехай вага ще й спадає пропорційно до $y^2$. Для нейрона з входами $\mathbf x$ і лінійним виходом $y=\mathbf w\cdot\mathbf x$ отримуємо
\begin{empheq}[box=\keybox]{equation}
\frac{\dd\mathbf w}{\dd t} \;=\; \eta\,y\,\bigl(\mathbf x - y\,\mathbf w\bigr) .
\label{eq:oja}
\end{empheq}
Правило, як і раніше, локальне: синапс знає свій вхід $x_i$, вихід $y$ і вагу $w_i$.

\begin{keyblock}
\begin{proposition}[Що знаходить правило Гебба]
\label{prop:oja}
Правило~\eqref{eq:oja} зводить вектор ваг до одиничної довжини вздовж напрямку, в якому входи змінюються найсильніше, --- до головного власного вектора матриці кореляцій входів $C=\langle\mathbf x\mathbf x^{\mathsf T}\rangle$~\cite{Oja1982}. Нейрон навчається видавати найвиразнішу одну величину, в яку можна стиснути його входи.
\end{proposition}
\end{keyblock}

Доведення. Усереднимо~\eqref{eq:oja} за входами: $\dd\mathbf w/\dd t=\eta\bigl(C\mathbf w-(\mathbf w^{\mathsf T}C\mathbf w)\,\mathbf w\bigr)$. Нерухомі точки --- власні вектори $C$ одиничної довжини. Якщо $\mathbf w$ близький до власного вектора $\mathbf e_k$ із числом $\lambda_k$, мала добавка вздовж іншого вектора $\mathbf e_j$ росте як $e^{\eta(\lambda_j-\lambda_k)t}$. Стійка лише та точка, для якої всі $\lambda_j<\lambda_k$, тобто головний вектор.

Є й варіант правила Гебба, який дивиться на час імпульсів. Якщо імпульс відправника прийшов за $s$ мілісекунд \emph{до} імпульсу отримувача, вага росте на $A_+e^{-s/\tau_+}$; якщо через $s$ мілісекунд \emph{після} --- зменшується на $A_-e^{-s/\tau_-}$, з $\tau_\pm$ порядку двадцяти мілісекунд. Таке правило виміряно на синапсах \nt{GLUT}-нейронів~\cite{BiPoo1998}. Воно посилює зв'язки, які \emph{могли бути причиною} відповіді, і послаблює ті, що запізнилися (рис.~\ref{fig:stdp}). Проженіть через мережу багато разів одну послідовність збуджень, і в ній самі виростуть зв'язки від кожної ланки до наступної --- ланцюжок із розділу~\ref{sec:glutcomputer}, зібраний без інженера.

\begin{figure}[t]
\centering
\begin{tikzpicture}[>=Stealth,x=0.07cm,y=1.3cm]
\draw[->,gray!70!black] (-66,0) -- (68,0) node[right,black] {\small $s$};
\draw[->,gray!70!black] (0,-1.2) -- (0,1.35) node[above,black] {\small зміна ваги};
\draw[very thick,blue!60!black] plot[domain=0.5:60,samples=50] (\x,{exp(-\x/20)});
\draw[very thick,red!70!black] plot[domain=-60:-0.5,samples=50] (\x,{-0.6*exp(\x/20)});
\draw[blue!60!black,thick] (0.5,0) -- (0.5,1);
\draw[red!70!black,thick] (-0.5,0) -- (-0.5,-0.6);
\foreach \x in {-40,-20,20,40}{\draw[gray!60!black] (\x,-0.04) -- (\x,0.04);}
\node[below,font=\scriptsize] at (20,-0.04) {$20\ms$};
\node[above,font=\scriptsize] at (-20,0.04) {$-20\ms$};
\node[align=left,font=\small,blue!60!black,anchor=west] at (22,0.95) {відправник раніше:\\ зв'язок посилюється};
\node[align=right,font=\small,red!70!black,anchor=east] at (-12,-0.95) {відправник пізніше:\\ зв'язок слабшає};
\end{tikzpicture}
\caption{Правило Гебба за часом. По горизонталі --- $s=t_{\text{отрим}}-t_{\text{відпр}}$, на скільки імпульс отримувача відстав від імпульсу відправника; $\tau_\pm=20\ms$, $A_-=0.6A_+$. Вікно завширшки кілька десятків мілісекунд --- того самого порядку, що й вікно збігу~\eqref{eq:window}.}
\label{fig:stdp}
\end{figure}

Але всі правила цього підрозділу вчать того, що \emph{часто}, а не того, що \emph{корисно}. Синапс, якому відомі лише $x$ і $y$, не знає, чи принесла дія сік, чи біль. Щоб навчатися на результаті, синапсу потрібен третій сигнал.

\section{Третій множник}

Третій сигнал приносить \nt{DOPA}: дофамінові нейрони розсилають широкомовно одне число --- помилку передбачення винагороди $\delta$~\eqref{eq:rpe} (розділ~\ref{sec:neurontypes}). Нехай зміна ваги пропорційна добутку трьох множників: активності відправника, активності отримувача і $\delta$.

Складність у тому, що винагорода надходить не тоді, коли мережа вибирала дію, а через сотні мілісекунд або секунди. На момент надходження $\delta$ добуток $xy$ давно дорівнює нулю, і синапсу потрібна пам'ять про те, що він брав участь. Найпростіша пам'ять --- знайоме нам RC-коло:
\begin{empheq}[box=\keybox]{equation}
\tau_e\,\frac{\dd e}{\dd t} \;=\; -e \;+\; x\,y ,
\qquad
\frac{\dd w}{\dd t} \;=\; \eta\,\delta(t)\,e(t) .
\label{eq:threefactor}
\end{empheq}
Величину $e$ називають \emph{слідом}~\cite{Fremaux2016}. Спільна активність лишає в синапсі мітку, яка тане за час $\tau_e$. Якщо за цей час надходить дофамін, вага змінюється зі знаком $\delta$; якщо ні, мітка зникає безслідно (рис.~\ref{fig:trace}). На \nt{GLUT}-синапсах виміряно: дофамін, що надійшов протягом секунди-двох після спільної активності, перетворює її на посилення зв'язку, а той, що надійшов пізніше, --- ні~\cite{Yagishita2014}. Вікна пластичності завдовжки в секунди знайшли й там, де третім сигналом слугує не дофамін, а сильне збудження самого отримувача~\cite{Bittner2017}.

\begin{figure}[t]
\centering
\begin{tikzpicture}[>=Stealth,x=7cm,y=1cm]
\fill[orange!12] (0.7,-0.2) rectangle (0.8,5.2);
% rows: x*y, delta, traces, weights
\foreach \y/\lab in {4.4/{$x\,y$},3.1/{$\delta$},1.4/{слід $e$},0/{вага $w$}}{
  \draw[gray!70!black] (0,\y) -- (1.42,\y);
  \node[left,font=\small] at (-0.02,\y+0.3) {\lab};}
\draw[very thick,blue!60!black] (0,4.4) -- (0,4.9) -- (0.2,4.9) -- (0.2,4.4);
\node[right,font=\scriptsize,blue!60!black] at (0.21,4.8) {відправник і отримувач активні разом};
\draw[very thick,orange!85!black] (0,3.1) -- (0.7,3.1) -- (0.7,3.7) -- (0.8,3.7) -- (0.8,3.1) -- (1.4,3.1);
\node[right,font=\scriptsize,orange!85!black] at (0.81,3.6) {надійшов дофамін};
% traces, each in units of its own maximum
\draw[very thick,blue!60!black] plot[domain=0:0.2,samples=20] (\x,{1.4+1.1*(1-exp(-\x))/(1-exp(-0.2))})
  plot[domain=0.2:1.4,samples=60] (\x,{1.4+1.1*exp(-(\x-0.2))});
\draw[thick,densely dashed,gray!50!black] plot[domain=0:0.2,samples=30] (\x,{1.4+1.1*(1-exp(-\x/0.05))/(1-exp(-4))})
  plot[domain=0.2:1.4,samples=80] (\x,{1.4+1.1*exp(-(\x-0.2)/0.05)});
\node[right,font=\scriptsize,blue!60!black] at (1.0,2.15) {$\tau_e=1\,\text{с}$};
\node[right,font=\scriptsize,gray!40!black] at (0.3,1.65) {$\tau_e=50\ms$};
% weights
\draw[very thick,blue!60!black] (0,0.25) -- (0.7,0.25) -- (0.8,0.8) -- (1.4,0.8) node[right,font=\scriptsize] {$\tau_e=1\,\text{с}$: вага зросла};
\draw[thick,densely dashed,gray!50!black] (0,0.2) -- (1.4,0.2) node[right,font=\scriptsize,gray!40!black] {$\tau_e=50\ms$: без змін};
% time axis marks
\foreach \x/\l in {0/0,0.2/0.2,0.7/0.7,1.4/1.4}{\draw[gray!60!black] (\x,-0.05) -- (\x,0.05) node[below=3pt,font=\scriptsize] {\l\,с};}
\draw[<->,thick,gray!50!black] (0.2,5.35) -- node[above,font=\scriptsize,black] {затримка винагороди $0.5\,\text{с}$} (0.7,5.35);
\end{tikzpicture}
\caption{Слід за правилом~\eqref{eq:threefactor}. Спільна активність триває $0.2\,\text{с}$, дофамін надходить через півсекунди після її кінця. Сліди показано в частках свого максимуму. Повільний слід ($\tau_e=1\,\text{с}$) на цей момент ще живий, швидкий ($\tau_e=50\ms$) давно розтанув.}
\label{fig:trace}
\end{figure}

\begin{keyblock}
\begin{proposition}[Навчання широкомовним сигналом]
\label{prop:threefactor}
Правило~\eqref{eq:threefactor} змінює лише ті синапси, які брали участь у роботі мережі за останні $\tau_e$, у бік, заданий спільним сигналом $\delta$. Широкомовний сигнал разом із локальним слідом дає адресну поправку. Затримка між дією і результатом допустима, доки вона не більша за $\tau_e$.
\end{proposition}
\end{keyblock}

\section{Чому це працює: шум як пошук}
\label{sec:dofagradient}

Чому правило~\eqref{eq:threefactor} веде до кращого? Відповідь дає шум із розділу~\ref{sec:code}. Нехай вихід нейрона $y=f(u)+\xi$, де $u=\sum_jw_jx_j$, а $\xi$ --- випадкове тремтіння з нульовим середнім і дисперсією $\sigma^2$. Винагорода $R$ залежить від $y$. Візьмімо як слід не просто $x_jy$, а його випадкову частину $x_j\xi$, а як третій множник --- помилку $\delta=R-\bar R$, де $\bar R$ --- очікувана винагорода. За малого шуму $R\approx R\bigl(f(u)\bigr)+R'\,\xi$, тому
\begin{empheq}[box=\keybox]{equation}
\bigl\langle \delta\,\xi\,x_j \bigr\rangle \;=\; \sigma^2\,R'\,x_j \;=\; \frac{\sigma^2}{f'(u)}\,\frac{\partial\langle R\rangle}{\partial w_j} .
\label{eq:perturbation}
\end{empheq}
Середній крок іде за градієнтом очікуваної винагороди~\cite{Williams1992}. Ніхто не обчислював похідних: мережа випадково відхилилася, дофамін повідомив, чи стало краще, і синапси, що брали участь, зсунулися у вигідний бік (рис.~\ref{fig:perturbation}). Шум, який у розділі~\ref{sec:code} заважав передавати повідомлення, тут стає пошуком.

\begin{figure}[t]
\centering
\begin{tikzpicture}[>=Stealth,x=2cm,y=3cm]
\draw[->,gray!70!black] (0,0) -- (4.2,0) node[below left,font=\small,black] {вихід $y$};
\draw[->,gray!70!black] (0,0) -- (0,1.2) node[left,font=\small,black] {винагорода $R$};
% reward hill
\draw[very thick,blue!60!black] plot[domain=0:4,samples=60] (\x,{max(0,1-((\x-3)/2.2)^2)});
\node[font=\scriptsize,blue!60!black,anchor=south] at (3,1.02) {найкращий вихід};
% noise around f(u)
\draw[thick,gray!60!black] plot[domain=0.6:2.4,samples=50] (\x,{0.28*exp(-((\x-1.5)/0.3)^2/2)});
\draw[densely dashed,gray!60!black] (1.5,0) -- (1.5,0.535);
\node[below,font=\small] at (1.5,0) {$f(u)$};
\node[font=\scriptsize,gray!50!black] at (1.5,-0.2) {тремтіння $\xi$};
% expected reward
\draw[densely dotted,gray!60!black] (0.5,0.535) node[left,font=\scriptsize,black] {$\bar R$} -- (2.1,0.535);
% two random deviations
\fill[green!45!black] (2.1,0.833) circle (0.035);
\draw[very thick,green!45!black] (2.1,0.535) -- (2.1,0.833);
\node[font=\scriptsize,green!45!black,anchor=west,align=left] at (2.18,0.6) {$\xi>0$: краще,\\ $\delta>0$};
\fill[red!70!black] (0.9,0.089) circle (0.035);
\draw[very thick,red!70!black] (0.9,0.535) -- (0.9,0.089);
\node[font=\scriptsize,red!70!black,anchor=east,align=right] at (0.83,0.27) {$\xi<0$: гірше,\\ $\delta<0$};
% resulting step
\draw[->,very thick,orange!85!black] (1.55,0.4) -- (1.95,0.4) node[right,font=\scriptsize,black] {крок};
\end{tikzpicture}
\caption{Шум як пошук~\eqref{eq:perturbation}. Вихід нейрона тремтить навколо $f(u)$. Випадкове відхилення праворуч (зелене) дало більше за очікуване, $\delta>0$; ліворуч (червоне) --- менше, $\delta<0$. В обох випадках добуток $\delta\,\xi$ додатний, і вага зсувається праворуч, туди, де винагорода росте. У середньому крок пропорційний нахилу $R'$: на вершині гірки відхилення в обидва боки однаково погані, і кроки гасять один одного.}
\label{fig:perturbation}
\end{figure}

\begin{keyblock}
\begin{proposition}[Спуск за градієнтом без зворотних проводів]
\label{prop:gradient}
Якщо в мережі є випадкові відхилення активності, а широкомовний сигнал дорівнює помилці передбачення винагороди і в кожній обстановці в середньому дорівнює нулю, правило~\eqref{eq:threefactor} у середньому змінює ваги за градієнтом очікуваної винагороди. Зворотні проводи й окрема фаза навчання для цього не потрібні.
\end{proposition}
\end{keyblock}

Тепер зрозуміло, чому дофамін повідомляє не про винагороду, а про \emph{помилку} її передбачення. Якби третім множником була сама винагорода $R\ge0$, середнє в~\eqref{eq:perturbation} не змінилося б, але з'явилися б дві біди. По-перше, до корисної частини додався б доданок $\bar R\,\xi x_j$: у середньому нуль, але сильний шум. По-друге, що гірше, справжній синапс не вміє відокремлювати в сліді $x_jy$ випадкову частину від невипадкової. За даних входів $x_jf(u)$ --- те саме число від спроби до спроби; якщо $\delta$ у цій обстановці в середньому дорівнює нулю, ця частина сліду нічого не змінює, а при $\delta\ge0$ вона раз у раз посилює все, що мережа робила, --- правильне і неправильне. Тому $\bar R$ має бути очікуванням \emph{у даній обстановці}, а не середнім за все життя. Обчислювати його --- справа критика, про якого нижче.

Ми перевірили це на моделі. Дві групи \nt{GLUT}-нейронів вибирають одну з двох дій через взаємне гальмування, як на рис.~\ref{fig:wta}; ваги від сигналу «почали» до груп спочатку рівні, а шум вирішує, хто переможе. Дія $A$ приносить сік з імовірністю $0.8$, дія $B$ --- з імовірністю $0.2$, і сік надходить через півсекунди після вибору. При $\tau_e=1\,\text{с}$ і $\delta=R-\bar R$ частка виборів $A$ за п'ятсот спроб зросла з $0.65$--$0.8$ у першій сотні до $0.96$--$1.0$ в останній, а ваги лишилися в межах $0.85$--$1.35$. При $\tau_e=50\ms$, меншому за затримку винагороди, мережа не навчилася нічого: частка виборів $A$ лишилася близько половини. При $\delta=R$ без віднімання очікування мережа теж стала вибирати $A$, але вага переможця за ті самі п'ятсот спроб зросла в тисячу разів і продовжувала рости; а при тому самому $\delta=R$ і вдесятеро більшій швидкості навчання мережа в одному з трьох прогонів назавжди закріпила свій перший випадковий вибір --- гірший.

\section{Ціна одного проводу}

Сигнал $\delta$ один на всіх, а шум, який він оцінює, --- сума тремтінь усіх $N$ нейронів, що впливають на результат. Кожен синапс бачить у $\delta$ свою корисну частину і чужий шум $N-1$ сусідів. Щоб виокремити своє, доводиться усереднювати за багатьма спробами, і час навчання росте пропорційно до $N$~\cite{Werfel2005}. Зворотне поширення такої плати не потребує.

\begin{keyblock}
\begin{proposition}[Ціна широкомовлення]
\label{prop:broadcastcost}
Час навчання за одним спільним сигналом росте пропорційно до кількості нейронів, чий шум впливає на результат. Таке навчання придатне для невеликих кіл, що вибирають із небагатьох варіантів, але не для налаштування мільярдів ваг.
\end{proposition}
\end{keyblock}

Мозок, схоже, так і розподіляє роботу: виходи \nt{DOPA}-нейронів ідуть насамперед до областей, які вибирають дії (розділ~\ref{sec:neurontypes}), а величезна \nt{GLUT}-мережа кори, де переважна більшість ваг, навчається здебільшого місцевими правилами, без зовнішнього вчителя. Раніше їх зводили до правил Гебба~\eqref{eq:oja}. Тепер дедалі більше даних свідчить, що кора має власну мету --- \emph{передбачати} свій наступний вхід, і тоді помилка обчислюється на місці, як різниця між передбаченим і тим, що надійшло~\cite{Keller2018}: учитель вбудований у саму мережу. Вікна пластичності завдовжки в секунди, де третім сигналом слугує сильне збудження самого отримувача~\cite{Bittner2017}, показують, що місцевий сигнал помилки в синапсів справді буває. Наскільки це загальне правило кори, --- гіпотеза.

\section{Звідки береться помилка: критик у неперервному часі}

Лишається питання, як самі \nt{DOPA}-нейрони обчислюють $\delta$. У програмах навчання з підкріпленням помилку передбачення рахують кроками $t,t+1,\dots$ В організмі кроків немає, тому запишімо її в неперервному часі~\cite{Doya2000}. Нехай $r(t)$ --- потік винагороди, а $V(t)$ --- очікування всієї майбутньої винагороди, в якій віддаленіша цінується менше:
\begin{equation}
V(t) \;=\; \int_t^{\infty} e^{-(s-t)/\tau_\gamma}\,r(s)\,\dd s .
\end{equation}
Продиференціювавши, отримаємо $\dd V/\dd t=V/\tau_\gamma-r$. Нев'язка цієї рівності і є помилкою передбачення:
\begin{empheq}[box=\keybox]{equation}
\delta(t) \;=\; r(t) \;+\; \frac{\dd V}{\dd t} \;-\; \frac{V}{\tau_\gamma} .
\label{eq:ctd}
\end{empheq}
Стала $\tau_\gamma$ --- горизонт планування, неперервний аналог коефіцієнта $\gamma$; за гіпотезою, його задає \nt{SERO} (підрозділ~\ref{sec:horizon}), і тоді~\eqref{eq:patience} --- правило, скільки варто чекати. Перевірмо три випадки (рис.~\ref{fig:ctdcases}). Засвітилася лампочка, що обіцяє сік: $V$ стрибком зросло, $\dd V/\dd t$ велике, сплеск. Надійшов обіцяний сік: $r>0$, але очікування в ту саму мить вичерпалося, $\dd V/\dd t\approx-r$, і сплеску немає. Сік не надійшов: очікування зникло без винагороди, $\dd V/\dd t<0$, провал.

\begin{figure}[t]
\centering
\begin{tikzpicture}[>=Stealth,x=1.15cm,y=1cm]
\foreach \col/\ttl/\lampon/\juice in {0/{несподіваний сік}/0/1, 1/{обіцяний сік}/1/1, 2/{сік не надійшов}/1/0}{
 \begin{scope}[xshift=\col*4.6cm]
  \node[font=\small] at (1.5,3.55) {\ttl};
  \foreach \yy in {2.4,1.2,0}{\draw[gray!60!black] (0,\yy) -- (3.1,\yy);}
  % reward r
  \ifnum\juice=1
    \fill[green!45!black] (1.95,2.4) rectangle (2.05,2.85);
    \pic[scale=0.55] at (2.0,3.1) {juice};
  \else
    \pic[scale=0.55] at (2.0,3.1) {nojuice};
  \fi
  % expectation V and error delta
  \ifnum\lampon=1
    \pic[scale=0.55] at (1.0,3.1) {lamp};
    \draw[very thick,blue!60!black] (0,1.2) -- (1,1.2) -- (1,1.45)
      plot[domain=1:2,samples=20] (\x,{1.2+0.25*exp((\x-1)/1.5)}) -- (2,{1.2+0.25*exp(1/1.5)}) -- (2,1.2) -- (3.1,1.2);
    \draw[very thick,red!70!black] (0,0) -- (0.95,0) -- (0.95,0.5) -- (1.05,0.5) -- (1.05,0) -- (3.1,0);
    \ifnum\juice=0
      \draw[very thick,red!70!black] (1.95,0) -- (1.95,-0.45) -- (2.05,-0.45) -- (2.05,0);
      \node[font=\scriptsize,red!70!black,anchor=west] at (2.1,-0.35) {провал};
    \else
      \node[font=\scriptsize,gray!50!black,anchor=north,align=center] at (2.0,-0.08) {$r$ і спад $V$\\ гасять одне одного};
    \fi
  \else
    \draw[very thick,blue!60!black] (0,1.2) -- (3.1,1.2);
    \draw[very thick,red!70!black] (0,0) -- (1.95,0) -- (1.95,0.5) -- (2.05,0.5) -- (2.05,0) -- (3.1,0);
  \fi
  \ifnum\col=0
    \node[font=\small,anchor=east] at (-0.1,2.55) {$r$};
    \node[font=\small,anchor=east] at (-0.1,1.35) {$V$};
    \node[font=\small,anchor=east] at (-0.1,0.1) {$\delta$};
  \fi
 \end{scope}}
\end{tikzpicture}
\caption{Три випадки для помилки~\eqref{eq:ctd}; згори донизу --- винагорода $r$, очікування $V$ і помилка $\delta$. Ліворуч очікування немає, і сік цілком --- несподіванка. Посередині лампочка стрибком піднімає $V$: це стрибок $\dd V/\dd t$, сплеск $\delta$ припадає на лампочку. Далі $V$ росте рівно так, як вимагає горизонт, і $\delta=0$; у момент соку винагорода $r$ і спад $V$ гасять одне одного. Праворуч $V$ падає без винагороди --- провал. Це ті самі сплеск, перенесення і провал, що на рис.~\ref{fig:rpe}, але тепер видно, звідки вони беруться.}
\label{fig:ctdcases}
\end{figure}

Що потрібно для~\eqref{eq:ctd} зі схем, які в нас уже є? Три речі.

\emph{Сама оцінка $V$.} Її видає група \nt{GLUT}-нейронів --- критик, --- ваги якої навчаються тим самим правилом~\eqref{eq:threefactor} з тим самим $\delta$: якщо $\delta>0$, очікування було занижене, і зв'язки, активні перед цим, посилюються.

\emph{Похідна $\dd V/\dd t$.} Її обчислює будь-яка схема, що реагує на зміни, а не на рівень: пряме збудження разом із затриманим гальмуванням через \nt{GABA}-посередника, як у детекторі напрямку з розділу~\ref{sec:gaba}, або синапс із виснаженням~\eqref{eq:depression} з розділу~\ref{sec:code}.

\emph{Знак в одному проводі.} Помилка буває і додатною, і від'ємною, а частота імпульсів --- лише додатною. Розв'язок той самий, що в \nt{SERO}-нейрона в розділі~\ref{sec:serocomputer}: \nt{DOPA}-нейрон --- ритмоводій. Його рівні кілька імпульсів на секунду означають $\delta=0$, пачка --- $\delta>0$, пауза --- $\delta<0$. Від'ємній помилці лишається менше місця: частоту не можна опустити нижче нуля. У справжніх \nt{DOPA}-нейронів провали справді мілкіші за сплески~\cite{Bayer2005}.

Що дофамін поводиться як $\delta$, виміряно. Як саме мозок обчислює $\dd V/\dd t$, --- гіпотеза.

\section{Актор і критик}

Зберімо все разом (рис.~\ref{fig:actorcritic}). Нейрони обстановки збуджують групи дій, що вибирають переможця через \nt{GABA} (твердження~\ref{prop:wta}), --- це \emph{актор}, --- і критика, що видає $V$~\cite{SuttonBarto2018}. \nt{DOPA}-нейрон отримує винагороду $r$ і $V$, обчислює~\eqref{eq:ctd} і розсилає $\delta$ усім синапсам актора і критика.

\begin{figure}[t]
\centering
\begin{tikzpicture}[>=Stealth,
  nrn/.style={circle,draw,very thick,fill=blue!6,minimum size=0.9cm,inner sep=0pt,font=\small},
  gab/.style={circle,draw,very thick,fill=red!10,minimum size=0.6cm,inner sep=0pt},
  dofa/.style={circle,draw,very thick,double,double distance=1.2pt,fill=orange!15,minimum size=1.35cm,inner sep=0pt,font=\scriptsize},
  inh/.style={-{Bar[width=7pt]},thick,red!70!black},
  bc/.style={->,thick,densely dashed,orange!85!black}]
\node[nrn] (S1) at (0,2.4) {$x_1$};
\node[nrn] (S2) at (0,1.2) {$x_2$};
\node[nrn] (S3) at (0,0) {$x_3$};
\node[left=0.15cm,align=right,font=\small] at (-0.5,1.2) {обстановка};
\node[nrn] (A) at (4,2.6) {$A$};
\node[nrn] (B) at (4,1.0) {$B$};
\node[gab] (G) at (5.2,1.8) {};
\draw[->,thick] (A) -- (G); \draw[->,thick] (B) -- (G);
\draw[inh] (G) to[bend left=35] (A);
\draw[inh] (G) to[bend right=35] (B);
\node[above,font=\small] at (4.6,3.05) {актор: вибір дії};
\node[nrn] (V) at (4,-1.1) {$V$};
\node[below,font=\small] at (4,-1.6) {критик};
\foreach \s in {S1,S2,S3}{\foreach \t in {A,B,V}{\draw[->,gray!60!black] (\s) -- (\t);}}
\node[dofa] (D) at (8.0,-1.1) {\nt{DOPA}};
\draw[->,thick] (V) -- node[above,font=\scriptsize] {$V,\ \dd V/\dd t$} (D);
\draw[->,thick,green!45!black] (10.0,-1.1) node[right,black,font=\small] {винагорода $r$} -- (D);
\pic[scale=1.1] at (10.6,-0.5) {juice};
\draw[bc] (D.north) -- (8.0,4.0) -- node[above,font=\small,black] {$\delta$ --- усім синапсам актора і критика} (2.2,4.0) -- (2.2,3.0);
\draw[bc] (D.south) -- (8.0,-2.4) -- (2.2,-2.4) -- (2.2,-0.6);
\end{tikzpicture}
\caption{Мережа, яка навчається вибирати дії. Сірі стрілки --- \nt{GLUT}-синапси зі змінними вагами; червоний кружок --- \nt{GABA}-нейрон; подвійний кружок --- \nt{DOPA}-ритмоводій, що обчислює $\delta$ за~\eqref{eq:ctd}; штрихові стрілки --- широкомовна розсилка $\delta$.}
\label{fig:actorcritic}
\end{figure}

Знак дії медіатора вибирає отримувач (твердження~\ref{prop:receptor}), і в області вибору дій частина нейронів несе дофамінові рецептори однієї родини, частина --- іншої. У дослідах на зрізах дофамін разом зі спільною активністю в перших посилює синапси, у других --- послаблює~\cite{Shen2008}; у моделі це означає, що в других знак $\delta$ у~\eqref{eq:threefactor} перевернутий. Перші навчаються того, що \emph{варто робити}, другі --- того, чого \emph{робити не варто}.

\section{Підсумок}

\begin{table}[t]
\centering
\footnotesize
\setlength{\tabcolsep}{4pt}
\renewcommand{\arraystretch}{1.15}
\begin{tabular}{>{\raggedright}p{2.6cm}>{\raggedright}p{2.7cm}>{\raggedright}p{3.1cm}>{\raggedright\arraybackslash}p{5.0cm}}
\toprule
Правило & Що знає синапс & Чого навчається мережа & Що відомо \\
\midrule
Гебб за частотами~\eqref{eq:oja} & $x$, $y$, свою вагу & стискання входів: головних напрямків & посилення за спільної активності виміряно; механізм обмеження ваг --- предмет досліджень \\
Гебб за часом & порядок імпульсів $x$ і $y$ у межах $\sim20\ms$ & причинних зв'язків, послідовностей & виміряно на \nt{GLUT}-синапсах \\
три множники~\eqref{eq:threefactor} & $x$, $y$, слід $e$, спільний сигнал $\delta$ & дій, які приносять винагороду & вплив дофаміну протягом секунд після активності виміряно; як головний механізм навчання вибору --- добре обґрунтована гіпотеза \\
критик~\eqref{eq:ctd} & те саме & очікування винагороди $V$ & схожість дофаміну з $\delta$ виміряно; схема, що обчислює $\dd V/\dd t$, --- гіпотеза \\
зворотне поширення помилки & персональну поправку для кожної ваги & усього, але потрібні зворотні проводи і такт & у мозку в цій формі не знайдено; шукають наближення до нього \\
\bottomrule
\end{tabular}
\caption{Правила навчання в мережі з \nt{GLUT}-, \nt{GABA}- і \nt{DOPA}-нейронів.}
\label{tab:learning}
\end{table}

Правила навчання зведено в табл.~\ref{tab:learning}. \nt{GLUT} несе дані, \nt{GABA} керує потоком, а \nt{DOPA} вирішує, які зміни в мережі залишити.

\section{Задачі}

\begin{exercise}[\lvlA]
\label{ex:06:1}
\emph{Скільки живе слід.} Спільна активність $xy=1$ триває $T_a=0.2\,\text{с}$, починаючи з $e=0$, а дофамін надходить через $d=0.5\,\text{с}$ після її кінця. (а)~У скільки разів слід~\eqref{eq:threefactor} у цей момент більший при $\tau_e=1\,\text{с}$, ніж при $\tau_e=50\ms$? (б)~За якого $\tau_e$ він найбільший?
\end{exercise}

\begin{exercise}[\lvlA]
\label{ex:06:2}
\emph{Дрейф правила Гебба.} Відправник і отримувач видають незалежні пуассонівські потоки по $10$ імпульсів на секунду, і кожна пара імпульсів змінює вагу за вікном рис.~\ref{fig:stdp} з $\tau_\pm=20\ms$, $A_-=0.6A_+$. На скільки $A_+$ зросте вага за годину цілком випадкової активності? За якого $\tau_-$ дрейф зникає?
\end{exercise}

\begin{exercise}[\lvlB]
\label{ex:06:3}
\emph{Кореляції, а не коваріації.} Правило твердження~\ref{prop:oja} знаходить головний власний вектор $C=\langle\mathbf x\mathbf x^{\mathsf T}\rangle$. Частоти додатні: $\mathbf x=\mathbf m+\mathbf n$, $\mathbf m=(\mu,0)$, коваріація шуму $\operatorname{diag}(s_1^2,s_2^2)$. За якої умови нейрон вивчить $\mathbf e_1$? Що він вивчить при $\mu=1$, $s_1=0.1$, $s_2=0.5$?
\end{exercise}

\begin{exercise}[\lvlA]
\label{ex:06:4}
\emph{Горизонт у сплеску.} Лампочка засвічується в непередбачуваний момент, сік розміру $R$ надходить через $L=1\,\text{с}$. Покажіть, що для вивченого очікування за~\eqref{eq:ctd} між лампочкою і соком $\delta\equiv0$, і знайдіть площу сплеску на лампочку при $\tau_\gamma=2\,\text{с}$ і $\tau_\gamma=4\,\text{с}$.
\end{exercise}

\begin{exercise}[\lvlB]
\label{ex:06:5}
\emph{Звідки береться ціна широкомовлення.} $N$ нейронів тремтять незалежно, $\xi_i\sim\mathcal N(0,\sigma^2)$, помилка $\delta=a\sum_i\xi_i$, синапс $j$ оцінює градієнт як $\delta\,\xi_jx_j$. Знайдіть відношення сигнал/шум однієї спроби. Один нейрон навчається за $100$ спроб по $5\,\text{с}$: скільки триватиме навчання при $N=10^4$ і $N=10^{10}$?
\end{exercise}

\begin{exercise}[\lvlB]
\label{ex:06:6}
\emph{Проєктування: схема, що обчислює $\delta$.} \nt{DOPA}-нейрон отримує $r$, $V$ з вагою $k_1$ і копію $V$, згладжену з $\tau_d$, з вагою $-k_2$. (а)~Доберіть $k_1$, $k_2$ так, щоб для повільних $V$ вихід був~\eqref{eq:ctd}. (б)~При $V=V_0e^{t/\tau_\gamma}$ істинна $\delta=0$. За якого $\tau_d$ хибна помилка схеми не перевищує $5\,\%$ від $V/\tau_\gamma$, якщо $\tau_\gamma=2\,\text{с}$?
\end{exercise}

\begin{exercise}[\lvlB]
\label{ex:06:7}
\emph{Моделювання: гранична затримка винагороди.} Додайте в \texttt{run()} копії \texttt{models/\allowbreak dofa\_learning.py} затримку винагороди $d$. Знайдіть частку виборів $A$ в останній сотні з $500$ спроб при $d=0.5\,\text{с}$ і $\tau_e=0.05$, $0.2$, $1$, $4\,\text{с}$ та при $\tau_e=1\,\text{с}$, $d=3\,\text{с}$. За якої частки сліду, що доживає до винагороди (задача~\ref{ex:06:1}), навчання зникає?
\end{exercise}

\begin{exercise}[\lvlC]
\label{ex:06:8}
\emph{Чи можна обійти ціну широкомовлення.} $N$ нейронів: $y_i=w_i+\xi_i$, $\xi_i\sim\mathcal N(0,\sigma^2)$, винагорода $R=-\sum_i(y_i-w_i^*)^2$, правило $\Delta w_i=\eta\,(R-\langle R\rangle)\,\xi_i$. За скільки спроб за найкращого $\eta$ помилка $\sum_i(w_i-w_i^*)^2$ спадає в $e$ разів? А якщо тремтять лише $m$ випадкових нейронів з $N$? Чи допомагає розріджений пошук?
\end{exercise}
